\documentclass[11pt,a4paper]{article}

\usepackage{geometry}
\usepackage{setspace}
\usepackage{parskip}
\usepackage{amsmath, amssymb}
\usepackage{booktabs}
\usepackage{multirow}
\usepackage{array}
\usepackage{graphicx}
\usepackage{caption}
\usepackage{xcolor}
\usepackage{hyperref}
\usepackage{listings}
\usepackage{xcolor}

% Define colors for syntax highlighting
\definecolor{codegreen}{rgb}{0,0.6,0}
\definecolor{codegray}{rgb}{0.5,0.5,0.5}
\definecolor{codepurple}{rgb}{0.58,0,0.82}
\definecolor{backcolour}{rgb}{0.95,0.95,0.92}

% Configure listings style
\lstdefinestyle{mystyle}{
    backgroundcolor=\color{backcolour},
    commentstyle=\color{codegreen},
    keywordstyle=\color{magenta},
    numberstyle=\tiny\color{codegray},
    stringstyle=\color{codepurple},
    basicstyle=\ttfamily\footnotesize\color{black},
    breakatwhitespace=false,
    breaklines=true,
    captionpos=b,
    keepspaces=true,
    numbers=left,
    numbersep=5pt,
    showspaces=false,
    showstringspaces=false,
    showtabs=false,
    tabsize=2
}
\lstset{style=mystyle}

\onehalfspacing

\title{\textbf{Dataset Prepration}\\
Multi-Task CNN for Mammogram Mass ROI Classification}
\author{Ali Saeidi Rad}
\date{27 February 2026}

\begin{document}

\maketitle

\begin{abstract}
    In this article i explain the dataset preapration for each dataset used.
\end{abstract}

\section{CBIS-DDSM}
\href{https://www.kaggle.com/datasets/awsaf49/cbis-ddsm-breast-cancer-image-dataset}{CBIS-DDSM} Dataset used for training Shape, Margin, Pathology, BI-RADS and Malignancy heads.
to prepare this dataset we first need to convert labels of Shape and Margin into standard form. existings labels are in Table.
because we aim for detecting mass, we only use the concatenated version of \texttt{mass\_case\_description\_train\_set.csv} and  \texttt{mass\_case\_description\_test\_set.csv} files.
To extract images and their corresponding masks we need to first extract the SeriesInstanceUID based on provided image file path.


\begin{lstlisting}[language=Python, caption=Extracting image file path]
jpeg = "path/to/jpeg/directory"
def SeriesInstanceUID(path):
    """Extract SeriesInstanceUID of the image file path"""
    return path.split("/")[-2]
def JPEGImage(series_uid):
    image_path = list(glob.glob(os.path.join(
        jpeg, str(series_uid),  "*.jpg"
    )))
    assert len(image_path) == 1
    return str(image_path[0])
df["ImageSeriesInstanceUID"] = df["image file path"].apply(
    SeriesInstanceUID
)
df["JPEGImage"] = df["ImageSeriesInstanceUID"].apply(JPEGImage)
\end{lstlisting}

\begin{lstlisting}[language=Python, caption=Extracting mask file path]
def StudyInstanceUID(path):
    return path.split("/")[-3]
def SeriesInstanceUID(path):
    return path.split("/")[-2]
def Name(path):
    return path["image file path"].split("/")[0] + \
    "_" + str(path["abnormality id"])

df["PatientName"] = df[
    ["image file path", "abnormality id"]
].apply(Name, axis=1)
df["MaskStudyInstanceUID"] = df[
    "ROI mask file path"
].apply(StudyInstanceUID)
df["MaskSeriesInstanceUID"] = df[
    "ROI mask file path"
].apply(SeriesInstanceUID)
df["PatientName"] = df[[
    "image file path", "abnormality id"]
].apply(Name, axis=1)
df["JPEGMask"] = df[
    [
        "PatientName",
        "MaskStudyInstanceUID",
        "MaskSeriesInstanceUID",
    ]
].apply(JPEGMask, axis=1)
\end{lstlisting}

after extracting image and mask file path, now we have to create a crop version of image,
based on the image and mask provided inside the dataset.

\begin{lstlisting}[language=Python, caption=Crop images based on provided mask]
output_dir = "desired/output/directory"
def JPEGCCrop(df):
    pname = df["PatientName"]
    image_path = df["JPEGImage"]
    image = cv2.imread(image_path, cv2.IMREAD_COLOR_RGB)
    mask = cv2.imread(df["JPEGMask"], cv2.IMREAD_GRAYSCALE)

    assert mask is not None
    assert image is not None

    _, thresh = cv2.threshold(mask, 127, 255, 0)
    contours, _ = cv2.findContours(
        thresh, 
        cv2.RETR_TREE, 
        cv2.CHAIN_APPROX_SIMPLE
    )
    list_of_pts = []
    for ctr in contours:
        list_of_pts += [pt[0] for pt in ctr]  # type: ignore
    center_pt = np.array(list_of_pts).mean(axis=0)
    clock_ang_dist = ClockwiseAngleDistance(center_pt)
    list_of_pts = sorted(list_of_pts, key=clock_ang_dist)
    ctr = np.array(list_of_pts).reshape((-1, 1, 2)).astype(np.int32)

    x, y, w, h = cv2.boundingRect(ctr)
    y_min, y_max = y, y + h
    x_min, x_max = x, x + w

    h, w = mask.shape
    pad_left, pad_bottom, pad_right, pad_top = 200, 200, 200, 200
    x_min = max(0, x_min - pad_left)
    y_min = max(0, y_min - pad_top)
    x_max = min(w, x_max + pad_right)
    y_max = min(h, y_max + pad_bottom)
    roi = image[y_min:y_max, x_min:x_max]

    crop_path = os.path.join(
        output_dir, 
        "images", 
        f"{pname}_{os.path.basename(image_path)}"
    )
    cv2.imwrite(crop_path, roi)

    return crop_path

df["JPEGCrop"] = df[
    ["JPEGImage", "JPEGMask", "PatientName"]
    ].apply(
    JPEGCCrop, axis=1
)
\end{lstlisting}

now we just need to convert labels of margin and shape into standard form.

\begin{lstlisting}[language=Python, caption=Labels conversations]
label_map = {
  "margin": {
    "SPICULATED": "SPICULATED",
    "Unknown": { "action": "drop" },
    "ILL_DEFINED": "ILL_DEFINED",
    "ILL_DEFINED-SPICULATED": "ILL_DEFINED",
    "MICROLOBULATED": "MICROLOBULATED",
    "MICROLOBULATED-ILL_DEFINED": "MICROLOBULATED",
    "MICROLOBULATED-SPICULATED": "MICROLOBULATED",
    "MICROLOBULATED-ILL_DEFINED-SPICULATED": "MICROLOBULATED",
    "OBSCURED": "OBSCURED",
    "OBSCURED-ILL_DEFINED": "OBSCURED",
    "OBSCURED-ILL_DEFINED-SPICULATED": "OBSCURED",
    "OBSCURED-SPICULATED": "OBSCURED",
    "OBSCURED-CIRCUMSCRIBED": "OBSCURED",
    "CIRCUMSCRIBED-ILL_DEFINED": "CIRCUMSCRIBED",
    "CIRCUMSCRIBED-OBSCURED": "CIRCUMSCRIBED",
    "CIRCUMSCRIBED-OBSCURED-ILL_DEFINED": "CIRCUMSCRIBED",
    "CIRCUMSCRIBED-MICROLOBULATED-ILL_DEFINED": "CIRCUMSCRIBED",
    "CIRCUMSCRIBED-MICROLOBULATED": "CIRCUMSCRIBED",
    "CIRCUMSCRIBED-SPICULATED": "CIRCUMSCRIBED",
    "CIRCUMSCRIBED": "CIRCUMSCRIBED"
  },
  "shape": {
    "ARCHITECTURAL_DISTORTION": { "action": "drop" },
    "LYMPH_NODE": { "action": "drop" },
    "FOCAL_ASYMMETRIC_DENSITY": { "action": "drop" },
    "ASYMMETRIC_BREAST_TISSUE": { "action": "drop" },
    "Unknown": { "action": "drop" },
    "ROUND": "ROUND",
    "ROUND-OVAL": "ROUND",
    "ROUND-LOBULATED": "ROUND",
    "ROUND-IRREGULAR-ARCHITECTURAL_DISTORTION": "ROUND",
    "IRREGULAR": "IRREGULAR",
    "IRREGULAR-ARCHITECTURAL_DISTORTION": "IRREGULAR",
    "IRREGULAR-FOCAL_ASYMMETRIC_DENSITY": "IRREGULAR",
    "IRREGULAR-ASYMMETRIC_BREAST_TISSUE": "IRREGULAR",
    "OVAL": "OVAL",
    "OVAL-LYMPH_NODE": "OVAL",
    "OVAL-LOBULATED": "OVAL",
    "LOBULATED": "OVAL",
    "LOBULATED-LYMPH_NODE": "OVAL",
    "LOBULATED-ARCHITECTURAL_DISTORTION": "OVAL",
    "LOBULATED-OVAL": "OVAL",
    "LOBULATED-IRREGULAR": "OVAL"
  },
  "pathology": {
    "MALIGNANT": "MALIGNANT",
    "BENIGN": "BENIGN",
    "BENIGN_WITHOUT_CALLBACK": "BENIGN"
  },
  "birads": {
    "0": "0",
    "1": { "action": "drop" },
    "2": "2",
    "3": "3",
    "4": "4",
    "5": "5"
  }
}
\end{lstlisting}

at the end we can split our data into train, validation and test.

\begin{lstlisting}[language=Python, caption=Train-val-test split based on subtlety provided inside the dataset]
def split_group(
    df: pd.DataFrame,
    train_frac: float,
    val_frac: float,
    test_frac: float,
    random_state=42,
    stratify: str = "pathology",
) -> Tuple[pd.DataFrame, pd.DataFrame, pd.DataFrame]:
    assert train_frac + val_frac + test_frac == 1.0

    train, temp = train_test_split(
        df,
        test_size=(1 - train_frac),
        stratify=df[stratify],
        random_state=random_state,
    )

    val, test = train_test_split(
        temp,
        test_size=test_frac / (val_frac + test_frac),
        stratify=temp[stratify],
        random_state=random_state,
    )

    return train, val, test
low = df[df["subtlety"] >= 4]
mid = df[df["subtlety"] == 3]
high = df[df["subtlety"] <= 2]

train_l, val_l, test_l = split_group(low, 0.6, 0.2, 0.2)
train_m, val_m, test_m = split_group(mid, 0.5, 0.25, 0.25)
train_h, val_h, test_h = split_group(high, 0.3, 0.3, 0.4)

train = pd.concat([train_h, train_m, train_l])
val = pd.concat([val_h, val_m, val_l])
test = pd.concat([test_h, test_m, test_l])
\end{lstlisting}

examples of images are in following figure.

\end{document}